% ETECOM https://ieee-etecom.org/
% Submission August 26
% Notification September 15
% Registration October 1
% Conference Octoer 26/27
% Track 1: AI, ML, Cybersecurity
% IEEE double-column format.
% 4-6 pages
\documentclass[conference]{IEEEtran}
\usepackage{cite}
\usepackage{amsmath,amssymb,amsfonts}
\usepackage{algorithmic}
\usepackage{graphicx}
\usepackage{textcomp}
\usepackage{xcolor}
\usepackage{float}
\usepackage{tikz}
\usetikzlibrary{positioning, shapes.geometric}
\usepackage{url}
\usepackage{subcaption}
\usepackage[compatibility=false]{caption}
\usepackage{array,tabularx,booktabs}
\usepackage{siunitx}
\usepackage{ragged2e}
\usepackage{etoolbox}
\usepackage{enumitem}
\usepackage{makecell}
\usepackage{multirow}
\usepackage{pifont}
\usepackage{listings} % For code blocks
\usepackage{comment}

% IEEE compliance fixes
\makeatletter
\patchcmd{\@maketitle}{\par}{\vspace*{-2ex}\par}{}{}
\makeatother

% Custom commands
\newcommand{\cmark}{\ding{51}}% ✓
\newcommand{\xmark}{\ding{55}}% ✗
\newcommand{\euro}{\text{\texteuro}}
\newcommand{\co}{CO\textsubscript{2}} % Proper CO₂ notation

% SI unit configuration
\sisetup{
 per-mode = symbol,
 group-digits = integer,
 output-decimal-marker = {.}
}

% Table formatting
\newcolumntype{Y}{>{\RaggedRight\arraybackslash}X}
\renewcommand{\arraystretch}{1.2}
\captionsetup[table]{font=small, labelfont=bf}
\setlist[itemize]{noitemsep, topsep=2pt}

% Code listing style (kept for potential future use if code is re-included, but not used in main text)
\definecolor{codegreen}{rgb}{0,0.6,0}
\definecolor{codegray}{rgb}{0.5,0.5,0.5}
\definecolor{codepurple}{rgb}{0.58,0,0.82}
\definecolor{backcolour}{rgb}{0.95,0.95,0.92}

\lstdefinestyle{mystyle}{
    backgroundcolor=\color{backcolour},
    commentstyle=\color{codegreen},
    keywordstyle=\color{magenta},
    numberstyle=\tiny\color{codegray},
    stringstyle=\color{codepurple},
    basicstyle=\ttfamily\footnotesize,
    breakatwhitespace=false,
    breaklines=true,
    captionpos=b,
    keepspaces=true,
    numbers=left,
    numbersep=5pt,
    showspaces=false,
    showstringspaces=false,
    showtabs=false,
    tabsize=2
}
\lstset{style=mystyle}

\begin{document}

\title{Beyond End-to-End Metrics: Stage-Oriented Evaluation of Cascaded Hybrid Intrusion Detection Systems}

\author{
\IEEEauthorblockN{Gina Bilinski\IEEEauthorrefmark{1} and Tianxiang Lu\IEEEauthorrefmark{1}}
\IEEEauthorblockA{\IEEEauthorrefmark{1}IU International University of Applied Science, Juri-Gagarin-Ring 152, 99084 Erfurt, Germany \\
Email: gina.bilinski@iu-study.org, tianxiang.lu@iu.org}
}
\maketitle

\begin{abstract}
Hybrid intrusion detection systems combine signature-based and anomaly-based techniques to detect both known and previously unseen attacks. While cascade-aware evaluation approaches have been explored in other machine learning domains, hybrid IDS studies typically evaluate complete architectures using overall performance measures such as Precision, Recall, or F$_1$-score, making it difficult to isolate the contribution of downstream detection stages. This paper presents a stage-oriented evaluation framework for cascaded hybrid IDS that explicitly quantifies the incremental contribution of downstream detection stages. The framework is instantiated using a pipeline that first analyzes network traffic with the signature-based intrusion detection system Snort. Subsequently, an Isolation Forest is applied exclusively to the residual traffic that remains unflagged by Snort. In addition to conventional classification metrics, the framework employs Conditional Detection Rate (CDR), Incremental Recall ($\Delta$Recall), and Incremental False Positive Rate ($\Delta$FPR) to
quantify  how many attacks missed by the first stage are recovered by the second stage and how the corresponding false-positive cost changes. On the CICIDS2017 dataset, the Isolation Forest achieves a CDR of 20.82\% on the Snort residuals, increasing the overall system Recall from 69.08\% to 75.52\%. This corresponds to a $\Delta$Recall of 6.44 percentage points and a $\Delta$FPR of 1.52 percentage points. The threshold analysis further shows that a higher recovery is possible at an increased false-positive cost, reaching a CDR of 48.98\%, a $\Delta$Recall of 15.14 percentage points and a $\Delta$FPR of 2.64 percentage points. Overall, the results demonstrate that residual-based anomaly detection can provide a measurable incremental contribution and that the proposed framework makes the contribution of individual detection stages explicit. The framework enables a more transparent evaluation and comparison of cascaded intrusion detection architectures.
\end{abstract}

\begin{IEEEkeywords}
Intrusion Detection System, 
Hybrid IDS,
Cascade Evaluation,
Anomaly Detection, 
Network Security, 
Evaluation Metrics
\end{IEEEkeywords}

\section{Introduction}
\label{sec:introduction}

Modern enterprise networks are continuously exposed to increasingly diverse and adaptive cyber attacks. Signature-based intrusion detection systems such as Snort remain effective for detecting known attack patterns, but they are limited to attacks for which suitable rules or signatures exist~\cite{roesch1999}. Anomaly-based intrusion detection can address this limitation by identifying deviations from normal network behavior. For this reason, hybrid intrusion detection systems combine both approaches in order to improve detection performance beyond a single detection method~\cite{maseno2022}.

In cascaded hybrid IDS architectures, detection components are applied sequentially. A signature-based detector can serve as the first stage, while a downstream anomaly detector analyzes the traffic that remains after this first stage. Although such systems are often evaluated through overall metrics such as Precision, Recall, Accuracy, or F$_1$-score, these metrics do not show how much the second stage contributes within the pipeline. For the evaluation of such a system, it is therefore necessary to examine how many attacks missed by the first stage are additionally detected and whether this detection gain is associated with an increase in false positives.

This paper addresses this evaluation gap by proposing a stage-oriented evaluation framework that quantifies the contribution of each stage independently. A Snort–Isolation Forest pipeline serves as a representative case study. Snort is used as the signature-based first stage. The traffic for which Snort generates no alert is treated as residual traffic and is then analyzed by an unsupervised Isolation Forest. The evaluation is conducted on the CICIDS2017 dataset and focuses on an evaluation methodology that enables a transparent assessment of individual detection stages within cascaded hybrid IDS architectures.

The contributions of this paper are as follows:
\begin{itemize}
    \item We propose a stage-oriented evaluation framework consisting of three complementary measures that quantify the effectiveness of downstream detection stages in cascaded IDS.
    \item We demonstrate the framework using the Isolation Forest specifically on the residuals of Snort, the network flows for which the signature-based stage generated no alert.
    \item We analyze how different anomaly thresholds influence the trade-off between additional attack detection and increased false positives on the CICIDS2017 dataset.
\end{itemize}

The remainder of this paper is organized as follows. Section~\ref{sec:related-work} reviews related work on hybrid intrusion detection systems and derives the research gap and research questions. Section~\ref{sec:methodology} introduces the experimental setup and evaluation methodology. Section~\ref{sec:result} presents the experimental results, followed by the discussion in Section~\ref{sec:discussion}. Finally, Section~\ref{sec:conclusion} concludes the paper.

%%%%%%%%%%%%%%%%%%%%%%%%%%%%%%%%%%%%%%%%%%%%%%%%%%%%%%%%%%%%%%%%%%%%%%%%%%%%%%

\section{Related Work}
\label{sec:related-work}

\subsection{Hybrid and Cascaded Intrusion Detection Systems}
\label{subsec:hybrid-cascaded-ids}

Hybrid intrusion detection systems have been widely studied, with
systematic reviews reporting improved overall detection performance
compared with individual approaches~\cite{maseno2022,abdulganiyu2023}.
However, evaluation commonly relies on aggregate metrics such as
Accuracy, Precision, Recall, F$_1$-score, or ROC-AUC, which provide
limited insight into the contribution of individual stages.

For Snort--machine-learning combinations, Bauskar et al.~\cite{bauskar2025}
use Snort logs as input to supervised models for further alert
classification, while Murgai et al.~\cite{murgai2024} transform Snort
rules into numerical features for unsupervised models including
Isolation Forest and K-Means. Although both approaches use Snort as an
input to downstream learning, neither specifically evaluates traffic
that remains undetected by Snort.

Multi-stage IDS architectures also process filtered traffic through
subsequent detection stages. Park and Lim~\cite{park2026}, for example,
apply later detection stages after an initial clustering-based filter
and report Precision and Recall for the anomaly-detection stage.
However, their first stage primarily filters normal traffic and thus
does not directly quantify the contribution of a downstream detector
after a signature-based attack detector. Sridharan et al.~\cite{sridharan2025}
use a structure closer to the present setting, forwarding traffic
without signature matches to a machine-learning anomaly detector.
Nevertheless, their evaluation focuses mainly on complete-system
performance and comparisons with standalone systems, leaving the
incremental contribution of the anomaly stage unclear.

Threshold selection represents a further limitation. Rookard and
Khojandi~\cite{rookard2025} determine an operating threshold by
optimizing G-Mean and Accuracy, but evaluate a single operating point
rather than how additional Recall and false-positive cost vary across
thresholds.

\subsection{Evaluation of Cascaded Machine Learning Systems}
\label{subsec:cascaded-ml-eval}

Beyond IDS, cascade evaluation explicitly considers the conditional
behavior of sequential stages. In the cascade architecture of Viola and
Jones~\cite{viola2001rapid}, increasingly complex classifiers operate
only on samples surviving previous stages, making stage-wise detection
and rejection behavior relevant to evaluation. Cascade generalization
similarly organizes classifiers sequentially, with later models
depending on preceding decisions~\cite{gamaCascadeGeneralization2000}.

This principle is directly relevant to cascaded hybrid IDS: a downstream
anomaly detector operates not on the original traffic distribution but
on residual traffic produced by preceding detection decisions.
Consequently, aggregate Recall or F$_1$-score cannot show how much
additional detection capability a downstream stage provides. While
conditional evaluation is established in cascaded machine learning,
IDS research has predominantly evaluated complete hybrid architectures
using aggregated metrics. The incremental contribution of anomaly
detectors operating on residual traffic therefore remains insufficiently
characterized.
%%%%%%%%%%%%%%%%%%%%%%%%%%%%%%%%%%%%%%%%%%%%%%%%%%%%%%%%%%%%%%%%%%%%%%%%%%%%%%

\subsection{Research Gap}
\label{subsec:research-gap}

Despite the prevalence of hybrid and multi-stage intrusion detection systems, existing studies largely report aggregate system-level metrics. This practice obscures the marginal impact and false-positive cost contributed by downstream stages. While stage-wise analysis is standard in the cascade learning literature~\cite{viola2001rapid,gamaCascadeGeneralization2000}, such evaluation is rarely applied within the IDS domain. Our work addresses this methodological gap by introducing a stage-oriented evaluation protocol for residual-forward cascades, enabling explicit quantification of each stage’s conditional and marginal contribution.
These observations motivate the central question of this paper: \textbf{How can the incremental contribution of downstream detection stages in cascaded hybrid intrusion detection systems be evaluated beyond conventional system-level metrics?} 

%%%%%%%%%%%%%%%%%%%%%%%%%%%%%%%%%%%%%%%%%%%%%%%%%%%%%%%%%%%%%%%%%%%%%%%%%%%%%%

\section{Methodology}
\label{sec:methodology}

\subsection{Stage-Oriented Evaluation Framework}
\label{subsec:framework}

Conventional IDS evaluation measures the performance of the complete detection
pipeline as a single classifier. While such end-to-end metrics are suitable for
comparing final detection performance, they do not reveal how individual stages
contribute to the overall system. In particular, the additional value of a
downstream detector in a cascaded IDS depends on the traffic that remains
undetected after previous stages.

Therefore, this work introduces stage-oriented evaluation as an evaluation
perspective for cascaded IDS architectures. Stage-oriented evaluation measures
the conditional contribution of each detection stage with respect to the
residual traffic produced by preceding stages. Instead of evaluating only the
final pipeline output, it explicitly quantifies the incremental detection
capability and the associated false-positive cost introduced by downstream
stages.

Formally, a cascaded IDS consists of sequential stages
$S_1, S_2, ..., S_n$. Each stage receives as input the residual set produced by
the previous stage:

\begin{equation}
R_i = X_i \setminus D_i
\end{equation}

where $X_i$ denotes the input traffic of stage $i$ and $D_i$ denotes the traffic
detected as malicious by that stage. The resulting residual set $R_i$
contains both benign flows that were not flagged by the previous stage and
malicious flows that remained undetected. The subset of
malicious flows within the residual set is referred to as the \emph{residual attack set}. This set represents the attack traffic available to downstream stages and defines the basis for measuring their conditional detection capability.

Consequently, the effectiveness of a
downstream stage is evaluated not on the original traffic distribution but on
the conditional distribution of residual traffic. The objective is therefore not to estimate the standalone detection capability of a stage, but its marginal contribution conditioned on the decisions of
previous stages.

This distinction is important because the residual composition directly depends
on preceding detection stages. A strong first-stage detector leaves fewer but
potentially harder-to-detect attacks, whereas a weaker first stage produces a
larger residual attack set. Therefore, the contribution of a downstream stage
cannot be interpreted independently from the preceding stages.

\subsection{Framework Instantiation}
\label{subsec:framework-instantiation}
A two-stage cascade consists of a first detector $S_1$ and a downstream
detector $S_2$. The first stage produces detected traffic $D_1$ and residual
traffic $R_1$. The second stage is evaluated only on $R_1$, while the final
pipeline decision is defined as the union of both stage outputs:

\begin{equation}
D_{Hybrid}=D_1 \cup D_2
\end{equation}

To demonstrate the applicability of the proposed evaluation perspective, a
two-stage IDS pipeline consisting of Snort and Isolation Forest is used as a
representative instantiation.

The first stage applies signature-based detection using Snort. Traffic without
a valid Snort detection forms the residual set:

\begin{equation}
R_{Snort}=X\setminus D_{Snort}
\end{equation}

This residual traffic contains both benign flows correctly ignored by Snort and
attack flows missed by the first stage. The second stage applies anomaly
detection on this residual set. The objective is therefore not to replace
Snort, but to quantify the additional detection capability gained by analyzing
the remaining unexplained traffic. Figure~\ref{fig:evaluation-logic} illustrates this instantiation. 


\begin{figure}[tb]
    \centering
    \includegraphics[width=0.45\textwidth]{images/evaluation_logic.png}
    \caption{Conceptual workflow of stage-oriented evaluation in a cascaded IDS.}
    \label{fig:evaluation-logic}
\end{figure}

%%%%%%%%%%%%%%%%%%%%%%%%%%%%%%%%%%%%%%%%%%%%%%%%%%%%%%%%%%%%%%%%%%%%%%%%%%%%%%

\subsection{Dataset and Preprocessing}
\label{subsec:dataset-preprocessing}
The framework is evaluated using CICIDS2017~\cite{sharafaldin2018} as a representative benchmark
dataset containing benign traffic and multiple attack scenarios.
%The CICIDS2017 dataset provides packet-level PCAP files as well as flow-based CSV files recorded over five consecutive days~\cite{sharafaldin2018}. 
This makes it suitable for the present experiment because Snort operates on packet captures, while the Isolation Forest and the final evaluation are performed on labeled network flows.

The Monday traffic of CICIDS2017 contains only benign traffic and is therefore used to train and calibrate the Isolation Forest. The traffic from Tuesday to Friday contains benign traffic and different attack classes and is used for the evaluation of the two-stage pipeline. 


Before applying the Isolation Forest, the data were cleaned by removing rows with missing or infinite values. Columns that contained only constant zero values and duplicate columns were removed because they do not contribute useful information to the anomaly detection model.

%%%%%%%%%%%%%%%%%%%%%%%%%%%%%%%%%%%%%%%%%%%%%%%%%%%%%%%%%%%%%%%%%%%%%%%%%%%%%%
\subsection{Experimental Protocol}
\label{subsec:experimental-protocol}
The evaluation follows three steps. First, Snort is executed on the packet
captures and the detected flows are removed from the original traffic,
creating the residual set. Second, the Isolation Forest is applied to this
residual set using thresholds calibrated on benign validation traffic. Third,
the outputs of both stages are combined to evaluate the final hybrid system.

\subsubsection{Stage 1: Signature-based Detection and Residual Construction}
\label{subsec:snort-residuals}
Snort identifies flows belonging to the first detection stage. The remaining
flows form the residual input of the second stage.
%%%%%%%%%%%%%%%%%%%%%%%%%%%%%%%%%%%%%%%%%%%%%%%%%%%%%%%%%%%%%%%%%%%%%%%%%%%%%%

\subsubsection{Stage 2: Residual Anomaly Detection}
\label{subsec:isolation-forest}
The Isolation Forest is trained on benign traffic and applied to the residual
flows generated by the first stage.


\subsection{Implementation Details}
\label{subsec:implementation-details}
\subsubsection{Flow-Level Alert Matching}
The official CICIDS2017 flow files provide timestamps only with minute-level
precision, whereas Snort alerts are recorded with microsecond precision. To
enable a temporally precise assignment of Snort detections to individual
network flows, the corresponding PCAP files were processed again with
CICFlowMeter, the tool originally used to generate the CICIDS2017 flow data~\cite{sharafaldin2018}.

The generated flow files were then combined with the labeled CICIDS2017 CSV files. The matching was performed using shared flow properties, including the five-tuple (source IP address, destination IP address, source port, destination port, and transport protocol), flow duration, and rounded start time. Additional packet and byte counts were used when needed to distinguish otherwise similar flows. This resulted in flow records with second-level timestamps and explicit end timestamps calculated from the flow duration.

Snort alerts were matched to flows using the timestamp and the five-tuple information. Since some detection rules identify attack behavior in server responses rather than the original client request, the matching was performed
bidirectionally. An alert was accepted only if its timestamp fell within the
time interval of the corresponding flow. For short flows below one second that
could not be matched using the exact interval, an additional tolerance window
of one second was evaluated. Ambiguous assignments were excluded to avoid
uncertain detections. Therefore, attacks detected by Snort but associated with
multiple possible flows remain in the residual set, resulting in a conservative
flow-level estimation of the first-stage performance.


Because the tolerance-based matching step has a noticeable influence on the number of matched flows, a sensitivity analysis with different tolerance windows was performed, while the definition of short flows remained unchanged.

\subsubsection{Detection Model Configuration}
The first detection stage was implemented using Snort 3 (version 3.12.2.0)
with the Talos LightSPD ruleset. The private network range of the CICIDS2017
victim network was configured as the protected internal network, and the
PortScan inspector was enabled because PortScan traffic represents one of the
attack categories contained in the dataset.

The second detection stage uses an Isolation Forest as an unsupervised anomaly detection model~\cite{liu2008} implemented with scikit-learn. The model was trained exclusively on benign Monday traffic to establish a representation of normal network behavior. The benign data were divided into training and validation subsets
using an 80/20 split. The model configuration used 300 trees.


The anomaly threshold was calibrated using the validation subset. In the main
configuration, the threshold was selected according to a target false-positive
rate of 1\% on benign validation traffic. Additional thresholds were evaluated
to analyze the trade-off between additional attack detection and false-positive
increase.

Since Isolation Forest relies on randomized tree construction, the evaluation
was repeated with ten different random states. The feature selection procedure,
training-validation split, and threshold calibration remained unchanged across
all repetitions.

\subsubsection{Feature Selection}
\label{subsubsec:feature-selection}
The Isolation Forest was first evaluated using all 67 numerical flow features available after preprocessing. Since flow-based IDS datasets often
contain redundant or highly correlated features, a feature reduction procedure
was applied to reduce redundancy and limit the influence of unstable features in the unsupervised model.
%improve the stability of the unsupervised model.

Feature selection was performed using only benign training traffic. Features
with no variation were removed first. Subsequently, features with a robust
coefficient of variation in the highest quartile were excluded. For feature
pairs with a correlation coefficient of $r>0.95$, the feature with the higher
robust coefficient of variation was removed to reduce redundant information.

From the remaining features, the 20 features with the lowest robust
coefficient of variation were selected. The selected features are listed in
Appendix~\ref{app:features} and mainly represent packet-length statistics,
header lengths, and inter-arrival-time characteristics.


%%%%%%%%%%%%%%%%%%%%%%%%%%%%%%%%%%%%%%%%%%%%%%%%%%%%%%%%%%%%%%%%%%%%%%%%%%%%%%

\subsection{Stage-Oriented Evaluation Metrics}
\label{subsec:evaluation-metrics}
The framework evaluates cascaded IDS pipelines from two complementary
perspectives. Conventional metrics quantify the final detection performance,
whereas stage-oriented metrics quantify the incremental contribution of
individual detection stages. 


The first metric is the Conditional Detection Rate (CDR). Unlike conventional Recall, CDR does not measure detection capability with respect to the original dataset, but with respect to the attack subset that survived previous detection stages. It measures the proportion of attacks missed by Snort that are recovered by the Isolation Forest. It is therefore calculated only on the attack flows contained in the Snort residuals. Here, $FN_{Stage1}$ denotes the number of attack flows missed by the first stage (e.g., Snort) and therefore correspond to the residual attack set, while $TP_{Stage2}$ denotes the number of these missed attack flows that are detected by the second stage (e.g., Isolation Forest). The CDR therefore corresponds to the Recall of the second stage on the attack part of the residual set. The separate term emphasizes that this value is conditional on the result of the first stage.

\begin{equation}
CDR = \frac{TP_{Stage2}}{FN_{Stage1}}
\label{eq:cdr}
\end{equation}


In addition, the change in Recall between the standalone Snort detection and the complete hybrid system is considered as Incremental Recall ($\Delta$Recall). $Recall_{Stage1}$ describes the proportion of all attacks detected by Snort, while $Recall_{Hybrid}$ describes the Recall of the complete two-stage system. $\Delta$Recall therefore shows by how many percentage points the Recall changes through the addition of the Isolation Forest.

\begin{equation}
\Delta Recall = Recall_{Hybrid} - Recall_{Stage1}
\label{eq:delta-recall}
\end{equation}

Since additional attack detections can be associated with additional false positives, the change in false-positive rate between both systems is also considered as Incremental FPR ($\Delta$FPR). $FPR_{Stage1}$ describes the false-positive rate of the standalone Snort detection, while $FPR_{Hybrid}$ describes the false-positive rate of the complete two-stage system. $\Delta$FPR therefore shows how much the proportion of benign network flows incorrectly classified as attacks changes through the second detection stage.

\begin{equation}
\Delta FPR = FPR_{Hybrid} - FPR_{Stage1}
\label{eq:delta-fpr}
\end{equation}

To analyze the influence of the Isolation Forest threshold, a $\Delta$Recall--$\Delta$FPR curve is additionally considered. The threshold determines the anomaly score above which a flow is classified as an attack and thereby directly controls the trade-off between additional attack detection and additional false alarms. For this purpose, a range of thresholds is evaluated and the resulting values of $\Delta$Recall and $\Delta$FPR are compared.

Together, CDR, $\Delta$Recall, and $\Delta$FPR provide complementary views of
the contribution of downstream detection stages. CDR quantifies the recovery
capability on the residual attack set, $\Delta$Recall measures the contribution
to the overall detection capability, and $\Delta$FPR captures the corresponding
false-positive cost. Evaluating these measures across different anomaly
thresholds enables a stage-oriented assessment of cascaded IDS pipelines beyond
conventional end-to-end performance metrics.

\subsection{Applicability of the Framework}
Although demonstrated using a two-stage Snort–Isolation Forest pipeline, the proposed evaluation methodology is applicable to arbitrary cascaded IDS architectures in which downstream stages operate on the residual output of preceding stages. Only the definition of the residual set and the detector outputs are required.
%%%%%%%%%%%%%%%%%%%%%%%%%%%%%%%%%%%%%%%%%%%%%%%%%%%%%%%%%%%%%%%%%%%%%%%%%%%%%%

\section{Results}
\label{sec:result}

\subsection{Residual Traffic Generated by the First Detection Stage}
\label{subsec:snort-results}
The first detection stage determines the residual traffic on which the
stage-oriented evaluation is performed. Snort generated 1,081,685 alerts, of
which 385,228 attack flows and 59,054 benign flows were matched to the
CICIDS2017 flow data, as shown in Figure~\ref{fig:snort-confusion}. Consequently, 172,410 attack flows and 1,683,583 benign
flows remained in the residual set and were passed to the second detection
stage.

\begin{figure}[htbp]
    \centering
    \includegraphics[width=0.50\textwidth]{images/confusion_matrix_snort.png}
    \caption{Confusion matrix of the Snort detection.}
    \label{fig:snort-confusion}
\end{figure}

The first stage achieves a Recall of 69.08\%, resulting in a residual attack
set of 172,410 flows. The corresponding Precision, F$_1$-score, and FPR are
86.71\%, 76.90\%, and 3.39\%, respectively.

The residual composition varies considerably between attack classes. PortScan
traffic is almost completely detected by the first stage, whereas DDoS and
several smaller attack classes remain largely in the residual set. Therefore,
the downstream detector is primarily evaluated on attack categories that are
challenging for signature-based detection. More details are provided in Appendix~\ref{app:detailed-results}. 

%%%%%%%%%%%%%%%%%%%%%%%%%%%%%%%%%%%%%%%%%%%%%%%%%%%%%%%%%%%%%%%%%%%%%%%%%%%%%%

\subsection{Stage-Oriented Evaluation of the Downstream Detector}
\label{subsec:if-results}
The downstream detector was evaluated on the residual set generated by Snort, consisting
of 172,410 missed attack flows and 1,683,583 benign flows (see confusion matrix in Figure~\ref{fig:if-confusion}). In this instantiation, the downstream detector is an Isolation Forest. At the main operating point with a target FPR of 1\%, the second stage detected 35,901 additional attacks while additionally misclassifying 26,421 benign residual flows. This corresponds to a Conditional Detection Rate (CDR) of 20.82\%. A detailed class-level breakdown is provided in Appendix~\ref{app:detailed-results}. The result quantifies the additional detection capability available from the residual attack set.

\begin{figure}[htbp]
    \centering
    \includegraphics[width=0.50\textwidth]{images/confusion_matrix_isolation_forest.png}
    \caption{Confusion matrix of the Isolation Forest on Snort residuals.}
    \label{fig:if-confusion}
\end{figure}

To assess the influence of feature selection, the Isolation Forest was compared
using all 67 numerical features and the selected 20-feature configuration.
The feature selection itself is performed only on benign training traffic and does not use attack labels. 
The reduced feature set improves Recall from 6.31\% to 20.82\% and provides
substantially lower variance across random states (0.62 percentage points versus
4.65 percentage points). Therefore, all subsequent stage-oriented evaluations use
the 20-feature configuration. 
%%%%%%%%%%%%%%%%%%%%%%%%%%%%%%%%%%%%%%%%%%%%%%%%%%%%%%%%%%%%%%%%%%%%%%%%%%%%%%

\subsection{Incremental Contribution of the Second Detection Stage}
\label{subsec:incremental-results}
The central objective of the stage-oriented evaluation is to quantify the
additional contribution of the downstream detector beyond the first stage.
The Isolation Forest recovers 35,901 of the 172,410 attacks remaining after
Snort. At the system level, this increases Recall
from 69.08\% to 75.52\%, corresponding to a $\Delta$Recall of 6.44 percentage
points. The additional detection capability is obtained at a cost of a
$\Delta$FPR of 1.52 percentage points.


\begin{table}[htbp]
\centering
\caption{Threshold analysis and impact on the hybrid IDS.}
\label{tab:threshold-analysis}
\scriptsize
\begin{tabular}{rrrrrrr}
\toprule
Target FPR (\%) & Threshold & CDR (\%) & $\Delta$Recall (pp) & $\Delta$FPR (pp) & Hybrid Recall (\%) & Hybrid FPR (\%) \\
\midrule
0.1 & 0.73 & 0.01 & 0.00 & 0.12 & 69.08 & 3.50 \\
0.5 & 0.69 & 0.45 & 0.14 & 0.89 & 69.22 & 4.28 \\
1.0 & 0.66 & 20.82 & 6.44 & 1.52 & 75.52 & 4.90 \\
2.0 & 0.62 & 48.98 & 15.14 & 2.64 & 84.23 & 6.03 \\
5.0 & 0.55 & 53.50 & 16.54 & 5.76 & 85.62 & 9.15 \\
10.0 & 0.50 & 53.86 & 16.65 & 10.14 & 85.74 & 13.53 \\
\bottomrule
\end{tabular}
\end{table}

The threshold analysis demonstrates the central trade-off captured by the
stage-oriented evaluation framework: the marginal attack recovery achieved by
the downstream stage versus the additional false-positive cost introduced into
the complete pipeline. Table~\ref{tab:threshold-analysis} shows how different Isolation Forest thresholds affect the additional benefit of the second detection stage. At target FPR values of 0.1\% and 0.5\%, the additional Recall remains close to zero. Between 1\% and 2\%, the CDR increases from 20.82\% to 48.98\%, while the $\Delta$Recall rises from 6.44 to 15.14 percentage points. 
This increase is mainly driven by the DDoS class, whose detection rate rises from 6.78\% to 53.33\% at the higher threshold setting (Appendix~\ref{app:detailed-results}). At the same time, the $\Delta$FPR increases from 1.52 to 2.64 percentage points. Higher target FPR values of 5\% and 10\% provide only a small additional Recall gain, but lead to a much stronger increase in false positives.

The $\Delta$Recall--$\Delta$FPR curve in Figure~\ref{fig:delta-recall-delta-fpr} was computed from 500 threshold values and illustrates how the additional detections change with the selected threshold. The blue points mark the comparison values from Table~\ref{tab:threshold-analysis}, while the red point marks the final threshold at a target FPR of 1\%. Starting from the 1\% target-FPR operating point, the curve rises steeply at a low $\Delta$FPR and reaches a $\Delta$Recall of approximately 16 percentage points, before it becomes increasingly flat.
 

\begin{figure}[htbp]
    \centering
\includegraphics[width=0.48\textwidth]{images/deltaRecall_deltaFPR.png}
    \caption{$\Delta$Recall--$\Delta$FPR curve of the Isolation Forest on Snort residuals.}
    \label{fig:delta-recall-delta-fpr}
\end{figure}

\subsection{Matching Sensitivity and Implementation Overhead}

The sensitivity analysis of the alert-flow matching in Table~\ref{tab:matching-tolerance} confirms that  temporal alignment has a substantial influence on the estimated first-stage performance. Increasing
the tolerance from 0.5 to 1.0 second increases Recall from 50.44\% to 69.08\%,
while further increases produce only marginal changes. Therefore, the selected
one-second tolerance provides a stable matching configuration.


\begin{table}[htbp]
\centering
\caption{Sensitivity analysis of the time tolerance used for alert-flow matching.}
\label{tab:matching-tolerance}
\scriptsize
\begin{tabular}{rrrrrrr}
\toprule
Tolerance (s) & Step 1 & Step 2 & TP & FP & Recall (\%) & FPR (\%) \\
\midrule
0.5 & 187,971 & 138,705 & 281,265 & 45,411 & 50.44 & 2.61 \\
1.0 & 187,971 & 256,311 & 385,228 & 59,054 & 69.08 & 3.39 \\
1.5 & 187,971 & 256,222 & 385,226 & 58,967 & 69.08 & 3.38 \\
2.0 & 187,971 & 256,788 & 385,782 & 58,977 & 69.18 & 3.38 \\
\bottomrule
\end{tabular}
\end{table}

The matching sensitivity analysis confirms the robustness of the first-stage
measurement. In addition, applying the downstream detector introduced only
approximately 1\% additional processing time compared with the complete
pipeline.

%%%%%%%%%%%%%%%%%%%%%%%%%%%%%%%%%%%%%%%%%%%%%%%%%%%%%%%%%%%%%%%%%%%%%%%%%%%%%%

\section{Discussion} 
\label{sec:discussion}
The results of the first detection stage show that a purely signature-based approach can already detect a large part of the attacks in the CICIDS2017 dataset. However, the aim of this work was not to optimize Snort itself, but to demonstrate the value of stage-oriented evaluation for quantifying the additional contribution of a downstream anomaly-based stage. With the freely available Talos LightSPD rules and without additional local rules, Snort reaches a Recall of 69.08\%.
The missing microsecond precision of the official CICIDS2017 flow data may have influenced the results. Although time corrections were applied, it cannot be fully excluded that individual alerts could not be assigned unambiguously to a flow despite a semantic relation. The sensitivity analysis of the matching tolerance shows that the selected tolerance of one second is methodologically plausible, since shorter tolerances clearly reduce the Snort Recall, while larger tolerances only lead to minor changes.

The Isolation Forest results show that the flows missed by Snort still contain patterns that can be used to identify additional attacks. The CDR of 20.82\% indicates what proportion of the attacks missed by Snort is subsequently detected by the Isolation Forest. This makes the contribution of the second stage visible without mixing it with the overall performance of the complete system. An important factor influencing the downstream detection performance is the
representation of residual traffic used by the anomaly detector. The comparison with the full 67-feature set shows that this selected representation improves the results and also makes them more stable.

The results for the complete hybrid system show that the second detection stage achieves a measurable increase in Recall with only a limited increase in FPR. The $\Delta$Recall of 6.44 percentage points directly describes how much the detection performance of the complete system improves through the Isolation Forest. As a result, the Recall of the hybrid system increases from 69.08\% to 75.52\%. At the same time, the $\Delta$FPR shows the additional false alarm rate associated with this gain. The FPR increases by 1.52 percentage points through the second stage. The additional detection performance is therefore associated with a smaller increase in false alarms than the corresponding increase in Recall. This makes it possible to evaluate the additional benefit of the second detection stage not only based on the additionally detected attacks, but also based on the additional false alarms it introduces. The robustness analysis further indicates that the measured additional benefit does not depend on a single random state of the Isolation Forest, since the derived measures vary only minimally across the ten runs.

The threshold analysis further illustrates that the incremental benefit of the second stage depends strongly on the selected threshold. With a target FPR of 2\%, the CDR increases to 48.98\%, meaning that almost half of the attacks missed by Snort are detected by the Isolation Forest. This illustrates the advantage of evaluating the downstream stage conditionally: the threshold decision can be interpreted directly in terms of recovered residual attacks and introduced false-positive cost. Compared with the 1\% setting, this provides a substantially higher additional detection gain, with $\Delta$Recall increasing from 6.44 to 15.14 percentage points, while the corresponding increase in $\Delta$FPR remains more limited, rising from 1.52 to 2.64 percentage points. The $\Delta$Recall--$\Delta$FPR curve confirms this relation. At the beginning of the curve, the $\Delta$Recall increases clearly, while the $\Delta$FPR increases only slightly. Afterwards, further threshold changes mainly lead to more additional false alarms without increasing the $\Delta$Recall to the same extent.

The results are consistent with previous observations that hybrid IDS architectures can improve detection performance compared with individual approaches~\cite{maseno2022}. For the first stage, a direct comparison of absolute values is only possible to a limited extent. Bauskar et al.~\cite{bauskar2025} report a standalone Snort detection rate of 81.1\% with an FPR of 16.3\%, but these values are not clearly attributed to a single dataset. Since the detection performance of signature-based systems depends on the deployed rule sets, identical datasets can yield strongly diverging results under different rules. The Recall of 69.08\% obtained here reflects the Talos LightSPD rules without attack-specific rule adaptations and is therefore only partially comparable with Snort results reported in the literature.

A quantitative comparison of the incremental contribution itself is hardly possible, since none of the reviewed studies reports measures equivalent to CDR or $\Delta$Recall. Even the structurally closest approach of Sridharan et al.~\cite{sridharan2025} evaluates only the complete system, and Murgai et al.~\cite{murgai2024} likewise report only aggregate metrics despite combining Snort with unsupervised models. From these studies it therefore cannot be derived how much of the reported detection gain actually originates from the second stage. This limitation motivates the need for evaluation approaches that separate the performance of the complete pipeline from the conditional contribution of individual detection stages.

The same applies to the threshold choice. Rookard and Khojandi~\cite{rookard2025} determine a single optimal threshold that maximizes G-Mean and Accuracy, whereas the present work examines the trade-off between additional Recall and additional false alarms across a range of target FPR values.

Several limitations should be considered. The entire evaluation is based on a single dataset, and Catillo et al.~\cite{catillo2022} show for CICIDS2017 that high detection performance on a public IDS dataset does not necessarily transfer to other network environments. The composition of the residuals further depends directly on the deployed rule set, so a different IDS or rule base would produce a different residual structure and correspondingly different results for the second stage. Consequently, stage-oriented metrics should not be interpreted as absolute detector quality measures, but as measurements of the added value of a stage within a specific cascade configuration. The same applies to the choice of the anomaly detector and the feature selection, since the Isolation Forest can only detect deviations that are reflected in the selected flow features. Finally, CDR, $\Delta$Recall, and $\Delta$FPR always describe the contribution of the second stage relative to the preceding first stage. When comparing different IDS architectures, they should therefore be interpreted together with the strength of the first detection stage and the composition of the residuals. 

%%%%%%%%%%%%%%%%%%%%%%%%%%%%%%%%%%%%%%%%%%%%%%%%%%%%%%%%%%%%%%%%%%%%%%%%%%%%%%

\section{Conclusion}
\label{sec:conclusion}
This paper introduced and evaluated a stage-oriented evaluation framework for
quantifying the incremental contribution of individual detection stages within
cascaded IDS pipelines. The framework was instantiated using an Isolation
Forest as a second detection stage in combination with Snort.
%This paper investigated the incremental contribution of an Isolation Forest as a second detection stage within a cascaded hybrid IDS pipeline. 
Instead of assessing only the overall performance of the complete system, the anomaly detector was evaluated specifically on the residuals of Snort, and its contribution was quantified through the derived measures CDR, $\Delta$Recall, and $\Delta$FPR.

On the CICIDS2017 dataset, Snort with the Talos LightSPD rules reaches a Recall of 69.08\%. The Isolation Forest recovers 20.82\% of the attacks missed by Snort, raising the Recall of the hybrid system to 75.52\% at a $\Delta$FPR of 1.52 percentage points. 
The threshold analysis further shows that the contribution of the downstream
stage depends strongly on the selected operating point, revealing the trade-off
between additional attack recovery and false-positive cost.
%The threshold analysis shows that a considerably higher contribution is possible at a moderately increased false alarm rate, with the CDR reaching 48.98\% at a target FPR of 2\%. The additional benefit of the second stage therefore depends strongly on the selected operating point.

Beyond these quantitative results, the main contribution of this work is methodological. Conventional system-level metrics show whether a hybrid IDS improves over a standalone detector, but they do not reveal which stage causes this improvement. A higher overall Recall can result from a strong first stage as well as from an actual contribution of the second stage. The proposed stage-oriented evaluation separates these effects and reports the additional detections and the additional false alarms of the downstream stage explicitly. This provides a transferable evaluation framework for assessing the contribution of individual stages in cascaded intrusion detection systems.

Future work could extend the approach to pipelines with more than two detection stages, where the contribution of each stage is evaluated relative to the
residual traffic produced by previous stages. If future studies report the contribution of downstream stages both on the residuals and relative to the complete system, different cascaded IDS architectures would become more directly comparable. Further research could also examine whether the observed contribution remains stable across other datasets, residual structures, and unsupervised detection methods, and what features are actually decisive for detection on residual traffic. 


% ---- Bibliography ----
\bibliographystyle{IEEEtran}
\bibliography{references}
\end{document}